\hypertarget{appendix-0976e86e5b9f}{}\label{referencestart:appendix-0976e86e5b9f}
\section{Off-axis stereo projection and disparity}\label{appendix-0976e86e5b9f-off-axis-stereo-projection-and-disparity}

\hypertarget{appendix-0976e86e5b9f-off-axis-stereo-projection-and-disparity}{}

\textbf{Curriculum depth:} Developed · \textbf{First seen:} Phase 11 · \textbf{Latest development:} Phase 11

\begin{allmightroledefinition}

\subsection{Definition and mechanism}\label{appendix-0976e86e5b9f-definition-and-mechanism}

\hypertarget{appendix-0976e86e5b9f-definition-and-mechanism}{}

Stereo cameras use two eye centres separated by a baseline, while each projection frustum is shifted so a chosen zero-parallax plane aligns in the images. Parallel eye orientations with asymmetric frusta avoid the vertical mismatch and discomfort often introduced by toe-in cameras. For rectified pinhole cameras, disparity is approximately \(\Delta x=f_xB/Z\), so inferred depth is \(Z=f_xB/\Delta x\) under the calibration and rectification assumptions. Depth uncertainty grows roughly with \(Z^2\) for fixed pixel disparity error. Distinguish interpupillary distance from lens and display-centre spacing; the rendered baseline is an optical/comfort design decision. Use consistent pixel units for focal length, baseline, and disparity before computing metric depth.

\end{allmightroledefinition}

\subsection{Prerequisites}\label{appendix-0976e86e5b9f-prerequisites}

\hypertarget{appendix-0976e86e5b9f-prerequisites}{}

\begin{itemize}
\tightlist
\item
  \href{../03-Geometry-Cameras-and-Image-Formation-P01/main.pdf\#appendix-c4e98b6e4001}{Stereo correspondence and epipolar geometry (Geometry, Cameras, and Image Formation, Part 1, Reference notes: Stereo correspondence and epipolar geometry)}
\item
  \href{../04-Light-Transport-and-Image-Synthesis-P01/main.pdf\#appendix-3ee9861904d5}{Camera calibration and bundle adjustment (Light Transport and Image Synthesis, Part 1, Reference notes: Camera calibration and bundle adjustment)}
\item
  \hyperref[appendix-cac01ebeeb11]{Axial depth and Euclidean range}
\end{itemize}

\subsection{Used by}\label{appendix-0976e86e5b9f-used-by}

\hypertarget{appendix-0976e86e5b9f-used-by}{}

\begin{itemize}
\tightlist
\item
  No downstream concept node is added in this stage.
\end{itemize}

\subsection{Full teaching explanation}\label{appendix-0976e86e5b9f-full-teaching-explanation}

\hypertarget{appendix-0976e86e5b9f-full-teaching-explanation}{}

\begin{itemize}
\tightlist
\item
  \href{../03-Geometry-Cameras-and-Image-Formation-P01/main.pdf\#02-foundations-ii-ch035-build-projection-and-stereo-geometry-around-real-camera-baselines}{Build projection and stereo geometry around real camera baselines (Geometry, Cameras, and Image Formation, Part 1, Ch. 4)}
\end{itemize}

\subsection{Source coverage}\label{appendix-0976e86e5b9f-source-coverage}

\hypertarget{appendix-0976e86e5b9f-source-coverage}{}

Supplied lesson parts: 1084--1088.

\subsection{Future keywords}\label{appendix-0976e86e5b9f-future-keywords}

\hypertarget{appendix-0976e86e5b9f-future-keywords}{}

zero-parallax plane, baseline, disparity uncertainty
